\documentclass[border=2pt]{standalone}
\usepackage{amsmath, amsthm, amsfonts}
\usepackage{tikz-cd, kotex}
\usepackage{quiver}
\usepackage{Operators}
\usepackage{palette}
\usepackage{diagram-style}
\begin{document}

%%FIG deformation of algebra (Deformation_Theory-1.svg)
\begin{tikzcd}
	C' \arrow[r, "\iota"] & C \\
	A' \arrow[r] \arrow[u] & A \arrow[u]
\end{tikzcd}

%%FIG isomorphism of deformations (Deformation_Theory-2.svg)
\begin{tikzcd}
	C' \arrow[rr, "\psi", "\sim"'] \arrow[rd, "\iota"'] & & C'' \arrow[ld, "\iota'"] \\
	& C &
\end{tikzcd}

%%FIG morphism of extensions (Deformation_Theory-3.svg)
\begin{tikzcd}
	0 \arrow[r] & M \arrow[r] \arrow[d, equal] & E_1 \arrow[r, "p_1"] \arrow[d, "\alpha"] & C \arrow[r] \arrow[d, equal] & 0 \\
	0 \arrow[r] & M \arrow[r] & E_2 \arrow[r, "p_2"'] & C \arrow[r] & 0
\end{tikzcd}

%%FIG trivial deformation of node (Deformation_Theory-4.svg)
\begin{tikzpicture}[line cap=round, line join=round]
  \path[use as bounding box] (-2,-2) rectangle (2,2);
  % Positive parameter: the horizontal component moves down;
  % the vertical component is common to every fibre.
  \foreach \offset/\tone in {0.16/accent3,0.36/accent2,0.60/accent1}{
    \draw[color=\tone,line width=0.8pt] (-2,-\offset) -- (2,-\offset);
  }
  \draw[black!50,line width=0.65pt] (-2,0) -- (2,0);
  \draw[black!50,line width=0.65pt] (0,-2) -- (0,2);
\end{tikzpicture}

%%FIG smoothing of node (Deformation_Theory-5.svg)
\begin{tikzpicture}[line cap=round, line join=round]
  \path[use as bounding box] (-2,-2) rectangle (2,2);
  \draw[black!50,line width=0.65pt] (-2,0) -- (2,0);
  \draw[black!50,line width=0.65pt] (0,-2) -- (0,2);
  % x=sqrt(t)*exp(s), y=sqrt(t)*exp(-s).
  % The finite range |s| <= log(2/sqrt(t)) keeps both coordinates
  % inside the square without evaluating t/x at or near x=0.
  \foreach \level/\tone in {0.015/accent3,0.085/accent2,0.25/accent1}{
    \pgfmathsetmacro{\smax}{ln(2/sqrt(\level))}
    \draw[color=\tone,line width=0.8pt]
      plot[domain=-\smax:\smax,variable=\s,samples=121]
      ({sqrt(\level)*exp(\s)},{sqrt(\level)*exp(-\s)});
    \draw[color=\tone,line width=0.8pt]
      plot[domain=-\smax:\smax,variable=\s,samples=121]
      ({-sqrt(\level)*exp(\s)},{-sqrt(\level)*exp(-\s)});
  }
\end{tikzpicture}

\end{document}
